% Parágrafo descritivo da categoria "Híbrido" usada nas figuras e tabelas
% de extração estruturada. Inclui referências cruzadas para
% \ref{fig:bubble-tecnica-ano}, \ref{fig:bubble-tecnica-subdominio} e
% \ref{tab:resumo-artigos}.

A categoria \emph{Híbrido}, que concentra \percHibrido{} dos artigos
incluídos no mapeamento (\totalIncluidos{} no total), agrupa trabalhos
cuja arquitetura predominante combina, de forma essencial, mais de
uma família técnica e que, por isso, não admitem classificação
unívoca em nenhuma das demais categorias. Quatro padrões recorrentes
foram identificados nesse grupo. O primeiro corresponde a
\textbf{empilhamentos neurais compostos}, nos quais camadas BiLSTM
e/ou CRF são acopladas sobre um \emph{transformer} pré-treinado
(e.g., RoBERTa--BiLSTM--CRF, BERT--BiLSTM--CRF). O segundo combina
\textbf{regras simbólicas com modelos estatísticos ou neurais},
tipicamente dicionários, listas controladas ou ontologias operando
em conjunto com classificadores supervisionados. O terceiro reúne
\textbf{pipelines multi-etapa}, nos quais o reconhecimento de
entidades é apenas um módulo de uma cadeia maior que inclui extração
de relações, sumarização ou construção de grafos de conhecimento. O
quarto agrupa \textbf{ensembles de aprendizado de máquina clássico}
(Random Forest, SGD, MLP) operando sobre representações textuais
como TF--IDF ou FastText, sem componente de aprendizado profundo.
Embora metodologicamente distintos, esses trabalhos compartilham a
propriedade de não serem redutíveis a uma única família de técnica,
motivo pelo qual aparecem consolidados sob o mesmo rótulo nas
Figuras~\ref{fig:bubble-tecnica-ano}
e~\ref{fig:bubble-tecnica-subdominio} e na
Tabela~\ref{tab:resumo-artigos}.
